\documentclass[12pt]{article}
\usepackage{amsmath, amssymb}
\usepackage[margin=1in]{geometry}
\usepackage{enumitem}

\setlength{\parindent}{0pt}
\setlength{\parskip}{6pt}

% Keep each problem, its working space, and its final answer together.
\newenvironment{workproblem}[1][4.2in]
    {\begin{minipage}[t][#1][t]{\linewidth}}
    {\end{minipage}}
\newcommand{\answerbox}[3]{%
    \par\noindent\hfill\fbox{%
        \begin{minipage}[t][#3][t]{#1}
            \small\textbf{#2}\par\vfill
        \end{minipage}}\par}
\newcommand{\setupanswer}{%
    \vfill\answerbox{\dimexpr\linewidth-2\fboxsep-2\fboxrule\relax}%
        {Final answer (integral setup)}{0.85in}}
\newcommand{\resultanswer}{%
    \vfill\answerbox{0.68\linewidth}{Final answer}{0.55in}}
\newcommand{\classificationbox}[1]{%
    \fbox{\begin{minipage}[t][0.46in][t]{0.29\linewidth}
        \small\textbf{#1}\par\vfill
    \end{minipage}}}
\newcommand{\classificationanswers}{%
    \par\vspace{0.08in}\noindent
    \classificationbox{Makes sense?}\hfill
    \classificationbox{Scalar / vector?}\hfill
    \classificationbox{Always zero?}\par\vspace{0.12in}}

\begin{document}

\textbf{MATH 261, Homework Set 9} \hfill Name: \underline{\hspace{6cm}}

\bigskip

%-----------------------------------------------------------------------

\begin{enumerate}[itemsep=0pt]
    \item \begin{workproblem}[3.85in]
    Let $S$ be the outward oriented sphere $x^{2}+y^{2}+z^{2}=1$, and let $\mathbf{F}=z\mathbf{i}-5\mathbf{k}$ be a vector field. Set up, but do NOT solve, the flux integral $\iint_{S} \mathbf{F} \cdot \mathbf{N} \, d\sigma$

    \setupanswer
    \end{workproblem}

    \item \begin{workproblem}[3.85in]
    Let $S$ be the surface obtained by rotating the curve $y=x^{3}$ for $1 \le x \le 3$ around the $x$-axis, oriented so the normal at the point $(2,8,0)$ is pointing in the positive $y$ direction. Let $\mathbf{F}=y\mathbf{i}+x\mathbf{j}+z\mathbf{k}$ be a vector field. Set up, but do NOT solve, the flux integral $\iint_{S} \mathbf{F} \cdot \mathbf{N} \, d\sigma$

    \setupanswer
    \end{workproblem}
    \newpage
    
    \item \begin{workproblem}
    Let $S$ be the surface given by the portion of the sphere $x^{2}+y^{2}+z^{2}=10$ which is above the plane $z=-1$ and below the plane $z=1$ oriented so the normal always points away from the origin. Let $\mathbf{F}=1\mathbf{i}+3\mathbf{j}+2\mathbf{k}$ be a vector field. Set up, but do NOT solve, the flux integral $\iint_{S} \mathbf{F} \cdot \mathbf{N} \, d\sigma$

    \setupanswer
    \end{workproblem}

    \item \begin{workproblem}
    Let $S$ be the surface given by the portion of the plane $2x+3y+z=6$ which lies in the first octant, oriented so that the normal always points in the positive $z$ direction. Let $\mathbf{F}=\langle 2,1,-2 \rangle$ be a vector field. Compute the flux integral $\iint_{S} \mathbf{F} \cdot \mathbf{N} \, d\sigma$

    \resultanswer
    \end{workproblem}
    \newpage

    \item \begin{workproblem}
    Let $S$ be the surface given by the portion of the plane $2x+3y+z=6$ which lies in the first octant, oriented so that the normal always points in the positive $z$ direction. Let $\mathbf{F}=\langle 1,0,-2 \rangle$ be a vector field. Compute the flux integral $\iint_{S} \mathbf{F} \cdot \mathbf{N} \, d\sigma$

    \resultanswer
    \end{workproblem}

    \item \begin{workproblem}
    Let $S$ be the surface given by the portion of the cone $z=\sqrt{x^{2}+y^{2}}$ which is above the plane $z=1$ and below the plane $z=2$ oriented so that the normal always points in the negative $z$ direction. Let $\mathbf{F}=\langle 1,1,z^{2} \rangle$ be a vector field. Compute the flux integral $\iint_{S} \mathbf{F} \cdot \mathbf{N} \, d\sigma$


    \resultanswer
    \end{workproblem}
    \newpage

    \item \begin{workproblem}
    Let $S$ be the outward oriented surface consisting of the upper half of the sphere with radius 1, $x^{2}+y^{2}+z^{2}=1$ for $z \ge 0$, and the disc $x^{2}+y^{2} \le 1$. Let $\mathbf{F}=\langle y,-x,1 \rangle$ be a vector field. Compute the flux integral $\iint_{S} \mathbf{F} \cdot \mathbf{N} \, d\sigma$

    \resultanswer
    \end{workproblem}

    \item \begin{workproblem}
    Let $S$ be the outward oriented surface consisting of the upper half of the sphere with radius 1, $x^{2}+y^{2}+z^{2}=1$ for $z \ge 0$, and the disc $x^{2}+y^{2} \le 1$. Let $\mathbf{F}=\langle 0,0,z \rangle$ be a vector field. Compute the flux integral $\iint_{S} \mathbf{F} \cdot \mathbf{N} \, d\sigma$

    \resultanswer
    \end{workproblem}
    \newpage

    \item \begin{workproblem}
    Let $S$ be the outward oriented surface consisting of the portion of the cylinder $x^{2}+y^{2}=1$ below the plane $z=1$ and above the plane $z=-1$, as well as the portions of planes $z=1$ and $z=-1$ which are within the cylinder (so the cylinder has a lid on both ends). Let $\mathbf{F}=\langle x,y,z^{2} \rangle$ be a vector field. Compute the flux integral $\iint_{S} \mathbf{F} \cdot \mathbf{N} \, d\sigma$

    \resultanswer
    \end{workproblem}

    \item \begin{workproblem}
    Let $\mathbf{F}=\langle y^{2}+z^{2},x^{2}+z^{2},x^{2}+y^{2} \rangle$ be a vector field, and let $C$ be the square bounded by the lines $x=\pm 1$, $y=\pm 1$, and $z=0$, traversed counterclockwise. Use Stokes' theorem to compute the line integral $\int_{C} \mathbf{F} \cdot \mathbf{T} \, ds$

    \resultanswer
    \end{workproblem}
    \newpage
     
    \item \begin{workproblem}
    Let $\mathbf{F}=\langle x^{2},2x,z^{2} \rangle$ be a vector field, and let $C$ be the ellipse $4x^{2}+y^{2}=4$, and $z=0$, traversed counterclockwise. Use Stokes' theorem to compute the line integral $\int_{C} \mathbf{F} \cdot \mathbf{T} \, ds$

    \resultanswer
    \end{workproblem}
    
    \item \begin{workproblem}
    Let $S$ be the outward oriented surface consisting of the top half of the sphere $x^{2}+y^{2}+z^{2}=16$ and the disc $x^{2}+y^{2} \le 16$ at height $z=0$. Let $\mathbf{F}=\langle x^{2},z^{2}y,zy^{2} \rangle$ be a vector field. Use Stokes' theorem to compute $\iint_{S} (\nabla \times \mathbf{F}) \cdot \mathbf{N} \, d\sigma$

    \resultanswer
    \end{workproblem}
    \newpage
    
    \item \begin{workproblem}
    Let $S$ be the upper half of the sphere $x^{2}+y^{2}+z^{2}=1$ for $z > 0$, oriented upward. Let $\mathbf{F}=\langle y,z,x \rangle$ be a vector field. Use Stokes' theorem to compute $\iint_{S} (\nabla \times \mathbf{F}) \cdot \mathbf{N} \, d\sigma$

    \resultanswer
    \end{workproblem}
    
    \item \begin{workproblem}
    Let $S$ be downward oriented surface consisting of the portion of the sphere $x^{2}+y^{2}+z^{2}=25$ which lies below the plane $z=-3$. Let $\mathbf{F}=\langle zy,-zx,e^{z^{2}} \rangle$ be a vector field. Use Stokes' theorem to compute $\iint_{S} (\nabla \times \mathbf{F}) \cdot \mathbf{N} \, d\sigma$

    \resultanswer
    \end{workproblem}
    \newpage

    \item \begin{workproblem}
    Let $S$ be the outward oriented sphere $x^{2}+y^{2}+z^{2}=1$. Let $\mathbf{F}=\langle z^{2},0,-xy \rangle$ be a vector field. Use Stokes' theorem to compute $\iint_{S} (\nabla \times \mathbf{F}) \cdot \mathbf{N} \, d\sigma$

    \resultanswer
    \end{workproblem}
     
    \item \begin{workproblem}
    Let $S$ be the downward oriented surface given by the portion of the cone $z=\sqrt{x^{2}+y^{2}}$ which is above the plane $z=1$ and below the plane $z=2$. Let $\mathbf{F}=\langle yz,0,z^{2} \rangle$ be a vector field. Use Stokes' theorem to compute $\iint_{S} (\nabla \times \mathbf{F}) \cdot \mathbf{N} \, d\sigma$

    \resultanswer
    \end{workproblem}
    \newpage

    \item \begin{workproblem}
    Let $S$ be the outward oriented surface consisting of the portion of the cylinder $x^{2}+y^{2}=1$ below the plane $z=1$ and above the plane $z=-1$. Let $\mathbf{F}=\langle y^{2},x,1 \rangle$ be a vector field. Use Stokes' theorem to compute $\iint_{S} (\nabla \times \mathbf{F}) \cdot \mathbf{N} \, d\sigma$

    \resultanswer
    \end{workproblem}

    \item \begin{workproblem}
    Let $S$ be the outward oriented surface consisting of the portion of the graph $z=10-x^{2}-y^{2}$ above the plane $z=1$. Let $\mathbf{F}=\langle y,0,x \rangle$ be a vector field. Use Stokes' theorem to compute $\iint_{S} (\nabla \times \mathbf{F}) \cdot \mathbf{N} \, d\sigma$

    \resultanswer
    \end{workproblem}
    \newpage

    \item (Conceptual problems) Let $f$ be a scalar function defined on $\mathbb{R}^{3}$ and let $\mathbf{F}$ be a vector field defined on $\mathbb{R}^{3}$. Assume that $f$ and all components of vector field $\mathbf{F}$ have continuous second partial derivatives. For each of the following expression, state (1) whether each expression makes sense mathematically. If it does, state (2) whether you get a scalar or vector field; and (3) whether it is always 0 or $\mathbf{0}$ if the object that it gives you is a vector.
    \begin{enumerate}
        \item \begin{minipage}[t]{\linewidth}
        $\nabla \times (\nabla \times \mathbf{F})$
        \classificationanswers
        \end{minipage}
        \item \begin{minipage}[t]{\linewidth}
        $\nabla \cdot (\nabla \cdot \mathbf{F})$
        \classificationanswers
        \end{minipage}
        \item \begin{minipage}[t]{\linewidth}
        $\nabla \times (\nabla \cdot \mathbf{F})$
        \classificationanswers
        \end{minipage}
        \item \begin{minipage}[t]{\linewidth}
        $\nabla \times (\nabla f)$
        \classificationanswers
        \end{minipage}
        \item \begin{minipage}[t]{\linewidth}
        $\nabla (\nabla \cdot \mathbf{F})$
        \classificationanswers
        \end{minipage}
        \item \begin{minipage}[t]{\linewidth}
        $\nabla \cdot (\nabla f)$
        \classificationanswers
        \end{minipage}
        \item \begin{minipage}[t]{\linewidth}
        $\nabla \cdot (\nabla \times \mathbf{F})$
        \classificationanswers
        \end{minipage}
    \end{enumerate}
\end{enumerate}

\end{document}
